\chapter{Introduction}\label{ch:introduction}
\chapterstatus{written}{2026-11-15}

\section{What FML is}

\fml{} is a library of building blocks, together with a few complete programs built from them. The building blocks
are templated C++ classes and functions in the namespace \nsp{FML}, one directory (\emph{module}) per topic. Most
modules are header-only; the ones with substantial compiled code (the Einstein--Boltzmann solver, recombination,
parts of the cosmology) have their own \file{Makefile}. Almost everything works both serially and in parallel with
MPI and OpenMP, with the same code: the parallel layout (\cref{ch:parallelism}) is part of the data structures, not
of the user code.

The complete programs are
\begin{itemize}
\item \textbf{SimulationCode} (\cref{ch:simulation-code}), a configurable particle--mesh / COLA / P$^3$M simulation code
      with a large library of cosmologies and theories of gravity, multiple particle species, on-the-fly analysis and
      initial conditions from a transfer-function file, from tabulated transfer tables or from the built-in
      Einstein--Boltzmann solver;
\item the \textbf{Einstein--Boltzmann solver} (\cref{ch:background,ch:recombination,ch:linear-perturbations,ch:cmb-spectra,ch:extended-models}),
      a line-of-sight code for the CMB and matter spectra in $\Lambda$CDM and some twenty extended models;
\item \textbf{SnapshotAnalyzer} (\cref{ch:snapshot-analyzer}), a standalone analysis program for simulation snapshots
      from any code;
\item \textbf{CosmologyFitting} (\cref{ch:sampling-fitting}), likelihoods and samplers built on the solver.
\end{itemize}

\section{Design principles}

A few principles recur throughout the library and explain many of the interfaces described in this book.

\paragraph*{One implementation of each physical quantity.} The expansion history, the quasi-static limits of
modified-gravity theories, the growth factors, the halo-model ingredients and the perturbation-theory kernels are each
implemented once and reused: the same \cls{Background} drives the Einstein--Boltzmann solver, the simulation time
stepping, the growth factors, spherical collapse and the halo model; the same \cls{QuasiStaticGravity} gives
$\mu(k,a)$ to the Boltzmann solver, to the LPT growth factors, to the one-loop hybrid bias spectra and to the
simulation code.

\paragraph*{Ownership is explicit.} Distributed data (particles in slabs, grids with local slices) and replicated data
(catalogues that every task holds in full) obey different rules, and each API states which it expects. Collective
functions must be called by all tasks.

\paragraph*{Units are explicit.} Box-unit positions in $[0,1)$ for particles and grids, $\hMpc$ and $\Mpch$ for
large-scale structure, SI or Mpc units in the Boltzmann code; \cref{ch:conventions} lists them and every chapter
states its own.

\paragraph*{Validation is part of a feature.} Every module has regression tests (analytic limits, manufactured
solutions, conservation laws, comparisons with external codes such as CLASS, CAMB, hi\_class, FAST-PT, velocileptors,
ReACT, ROCKSTAR or FINUFFT). The validation boxes in this book summarise them; the scope statements of the module
guides say what a passed test does and does not establish.

\section{Building}

\fml{} needs a C++17 compiler. The optional dependencies are switched on with make variables: \code{USE_MPI},
\code{USE_OMP}, \code{USE_FFTW} (FFTW3, with its MPI build for MPI runs), \code{USE_FFTW_THREADS}, \code{USE_GSL},
\code{USE_LUA} (parameter files), \code{USE_CGAL} (triangulations), \code{USE_HDF5} and \code{USE_PYTHON}. Machine
specific paths go in \file{FML/Makefile.local}; the shared rules are in \file{FML/Makefile.config}.
\begin{paramfile}[Makefile.local]
CXX_SERIAL   = c++
CXX_MPI      = mpicxx
LOCAL_PREFIX = /opt/homebrew      # where FFTW, GSL, Lua, ... are installed
\end{paramfile}
The test suite is built and run with \code{make -C FML test} (optionally with \code{MPIRUN=mpirun NP_LIST="1 2 4"}),
the examples with \code{make -C FML examples}, the simulation code with \code{make -C FML/SimulationCode} (or its CMake
build) and the Einstein--Boltzmann solver with \code{make} in \file{Cosmology/EinsteinBoltzmannSolver}. The file
\file{FML/BUILDING.md} is the reference for all build options.

\section{Map of the library}

\Cref{tab:module-map} groups the modules by the parts of this book. The full mapping of directories to sections is in
\cref{app:module-index}.

\begin{table}[htbp]
\centering\small
\caption{The modules of \fml{} grouped by the parts of this book.}\label{tab:module-map}
\begin{tabularx}{\textwidth}{@{}lX@{}}
\toprule
Part & Modules \\
\midrule
Infrastructure & \file{Global}, \file{PencilDomain}, \file{PeanoHilbertDomain}, \file{FFTWGrid}, \file{MPIGrid},
  \file{AdaptiveMPIGrid}, \file{MPIParticles}, \file{ParticleTypes}, \file{ParticlesInBoxes}, \file{KDTree},
  \file{Interpolation}, \file{Smoothing}, \file{NUFFT}, \file{Math}, \file{ODESolver}, \file{Spline}, \file{FFTLog},
  \file{RandomGenerator}, \file{LatinHyperCube}, \file{MultigridSolver}, \file{GadgetUtils}, \file{RamsesUtils},
  \file{FileUtils}, \file{ParameterMap}, \file{LuaFileParser}, \file{Units}, \file{Timing}, \file{MemoryLogging},
  \file{Performance}, \file{Plotting}, \file{Visualization}, \file{Chunkpix}, \file{ExampleData},
  \file{MakePythonWrapper}, \file{PythonAnalysis} \\
Linear theory & \file{Cosmology/Background}, \file{Cosmology/Recombination}, \file{Cosmology/Perturbations},
  \file{Cosmology/Spectra}, \file{Cosmology/EinsteinBoltzmannSolver}, \file{Cosmology/Growth},
  \file{TransferProviders}, \file{InitialConditions} \\
Nonlinear models & \file{PerturbationTheory}, \file{SphericalCollapse}, \file{HaloModel}, \file{AxionHMcode},
  \file{Emulator} \\
Initial conditions & \file{RandomFields}, \file{LPT}, \file{GlassIC} \\
Simulations & \file{NBody}, \file{ShortRangeGravity}, \file{SimulationCode}, \file{AMRCode}, \file{Hydrodynamics}, \file{SPH},
  \file{WaveDarkMatter}, \file{PGD}, \file{Baryonification} \\
Structures & \file{FriendsOfFriends}, \file{PhaseSpaceHaloFinder}, \file{HaloHistory}, \file{CosmicWeb},
  \file{TidalTensor}, \file{Morphology}, \file{VoidFinders}, \file{Triangulation} \\
Statistics & \file{ComputePowerSpectra}, \file{PairCounting}, \file{ThreePointCorrelation}, \file{CountsInCells},
  \file{Environment}, \file{ConnectedFourPoint}, \file{EnsembleStatistics}, \file{HOD} \\
Surveys and inference & \file{Survey}, \file{WeakLensing}, \file{FieldLevel}, \file{MCMC}, \file{CosmologyFitting} \\
Applications & \file{Analysis}, \file{SnapshotAnalyzer} \\
\bottomrule
\end{tabularx}
\end{table}
